\yibinopeningchapter

结构化排版系统通过源文件描述章节、引用和交叉关系，再由编译器生成稳定版面。
LaTeX 的设计强调内容组织与排版命令的组合\yibincite{lamport1994latex}；DOCX 则以 Office Open XML
文档结构保存段落、样式和分节信息\yibincite{iso29500-1-2016}。双格式输出因此不能只依赖文件扩展名转换，
还需要明确语义映射和版式校验边界。

\section{研究问题}

本案例关注两个可验证的问题：外部论文源能否在不复制到模板仓库的情况下完成构建；
PDF 与 DOCX 输出中的标题、图表、公式、注释和参考文献能否保持可检查的一致关系。

\section{方法边界}

Pandoc 能够在多种标记格式之间进行语义转换\yibincite{pandoc-manual}，但复杂分页、动态域和学校专用封面仍需由目标格式的专门处理逻辑完成。
因此，本案例只验证受控命令子集，不把 PDF 与 Word 的逐页像素一致作为目标。
